\documentclass[11pt]{article}
\usepackage{amsmath,amssymb,amsthm,hyperref}
\newtheorem{theorem}{Theorem}
\newtheorem{lemma}[theorem]{Lemma}
\newtheorem{corollary}[theorem]{Corollary}
\newcommand{\E}{\mathbb E}
\newcommand{\Prob}{\mathbb P}
\newcommand{\Var}{\operatorname{Var}}
\title{A matching power lower bound for random palindrome blocks}
\author{Research note; independently reviewed by two other agents}
\date{30 September 2026}
\begin{document}
\maketitle

Let $m$ independent uniformly relabelled palindrome blocks on $n$ labels
be concatenated, and put $R_\sigma=n!P(\sigma)$.
All constants below are absolute; inequalities stated for sufficiently
large $n$ can be absorbed into one absolute threshold $n_0$.

\begin{theorem}\label{thm:lower}
There exist constants $c,C>0$ and $n_0$ such that, whenever
$n\geq n_0$ and $m\geq Cn^{13/10}$,
\[
 \Prob\left(\max_{\sigma\in S_n}|R_\sigma-1|
       \geq c\frac{n^{13/4}}{m^{5/2}}\right)\geq1-\frac Cn.
\]
This is a lower bound for the specified random construction, not for
arbitrary collections of dice. Together with the accompanying uniform
upper bound, it matches the powers of $n$ and of the error tolerance,
up to the factor $(\log n)^{1/5}$ in the sufficient face count.
\end{theorem}

\section{Triple projection and central binomial weights}
Represent each block permutation by independent continuous priorities
$Y_{i,1},\ldots,Y_{i,n}$. For a target order, a consecutive triple has
internal ratio $3!q$ equal to zero if its middle priority is the smallest,
and equal to $3/2$ otherwise. Put
\[
 \delta_{i,a}=\frac12-\frac32\mathbf1\{\text{the middle of the triple
 starting at }a\text{ has smallest priority}\}.
\]
It is centered. The part of the exact Hoeffding expansion using one block
and precisely three labels is
\begin{equation}\label{eq:triple}
 S_\sigma=h\sum_{i=1}^m\sum_{a=0}^{d}w_i(a)\delta_{i,a},
 \quad d=n-3,\quad
 w_i(a)=\binom da p_i^a(1-p_i)^{d-a},\quad
 p_i=\frac{i-1}{m-1},
\end{equation}
where
\[
 h=\binom n3m^{-3}(1-1/m)^{n-3}.
\]
For $m\geq n\geq6$, $h\asymp n^3/m^3$. Formula~\eqref{eq:triple}
follows either by conditioning on the three labels occupying a specified
block, or by setting $\ell=1,k_1=3$ in the exact gap-count coefficient
formula of the companion upper-bound note.

\begin{lemma}\label{lem:binomial}
For sufficiently large $d$, integers $a$ with $d/4\leq a\leq3d/4$, and
$m\geq d+3$,
\[
 \sup_{0\leq p\leq1}\binom da p^a(1-p)^{d-a}\leq\frac C{\sqrt d},
 \qquad
 c\frac m{d^{3/2}}\leq\sum_{i=1}^m w_i(a)^2
                \leq C\frac m{d^{3/2}}.
\]
\end{lemma}
\begin{proof}
The maximum is attained at $p=a/d$. Stirling's factorial bounds, uniformly
for $a/d\in[1/4,3/4]$, give the first estimate. The beta integral is
\[
 \int_0^1\binom da^2p^{2a}(1-p)^{2d-2a}\,dp
 =\frac{\binom da^2}{(2d+1)\binom{2d}{2a}}
 \asymp d^{-3/2},
\]
again uniformly by the same bounds. The integrand is nonnegative and
unimodal with maximum $O(1/d)$. Its sum over the equally spaced grid with
$m$ points differs from $(m-1)$ times its integral by at most twice that
maximum. Since $m\geq d+3$, this error is a vanishing proportion of the
main term as $d\to\infty$.
\end{proof}

\section{Independent local swaps produce a large range}
Start from the target order $(1,\ldots,n)$, using positions numbered from
zero. Choose $J\asymp n$ disjoint windows of six consecutive positions in
the middle third of this order. In each window allow the two middle
positions to be swapped, independently of the choices in the other
windows. This defines a family $\mathcal F$ of $2^J$ target orders.
Every affected triple start lies in $[d/4,3d/4]$ for sufficiently large $n$.

If the two swapped positions of window $j$ are $u,u+1$, only the four
triple starts $u-2,u-1,u,u+1$ are affected. Define $D_j$ to be the change
in $S_\sigma$ caused by making this swap, leaving all other windows in
the original order. Its form is
\[
 D_j=h\sum_{i=1}^m B_{i,j},\qquad
 B_{i,j}=\sum_{a=u-2}^{u+1}w_i(a)
              (\delta_{i,a}^{\mathrm{swapped}}-\delta_{i,a}^{\mathrm{original}}).
\]
The $B_{i,j}$ are centered and independent over both $i$ and $j$: each
uses only the six priorities in its own window of its own block.
Moreover the effect of the choices is additive, so the range of $S$ on
the family is exactly
\begin{equation}\label{eq:range}
 \max_{\sigma\in\mathcal F}S_\sigma-
 \min_{\sigma\in\mathcal F}S_\sigma=\sum_{j=1}^J|D_j|.
\end{equation}

A simple event prevents cancellation in the variance of each local swap.
Suppose the relative ranks of the six priorities are
\[
 (6,1,3,2,4,5).
\]
Swapping the middle positions changes them to $(6,1,2,3,4,5)$ and removes
exactly the valley with triple start $u$. The other three affected valley
indicators are unchanged. Thus $B_{i,j}=(3/2)w_i(u)$ on this event.
The reverse rank event gives $B_{i,j}=-(3/2)w_i(u)$. Each event has
probability $1/6!$. Consequently
\begin{equation}\label{eq:variance}
 \E B_{i,j}^2\geq\frac1{160}w_i(u)^2.
\end{equation}
On the other hand, each valley-indicator difference is at most one in
absolute value, whence
\[
 \E B_{i,j}^2\leq9\sum_{a=u-2}^{u+1}w_i(a)^2,
 \qquad |B_{i,j}|\leq Cn^{-1/2}
\]
by Lemma~\ref{lem:binomial}. Independence over blocks and the lemma give,
with $V=h^2m/n^{3/2}$,
\[
 cV\leq\E D_j^2\leq CV.
\]
The fourth-moment formula for a sum of independent centered variables and
the uniform bound on $B_{i,j}$ imply
\[
 \E D_j^4\leq3(\E D_j^2)^2
       +\frac{Ch^2}{n}\E D_j^2
 \leq CV^2.
\]
The last step uses $(h^2/n)/V=\sqrt n/m\leq1$.
Moment interpolation now yields
\[
 \E|D_j|\geq\frac{(\E D_j^2)^{3/2}}{(\E D_j^4)^{1/2}}
               \geq c\sqrt V.
\]
The different $D_j$ are independent. Therefore the sum in
\eqref{eq:range} has expectation at least $cJ\sqrt V$ and variance at most
$CJV$. Chebyshev's inequality shows that, with probability $1-O(1/n)$,
\begin{equation}\label{eq:triplelower}
 \max_{\sigma\in\mathcal F}|S_\sigma|
 \geq\frac12\sum_j|D_j|
 \geq c n h\sqrt{m/n^{3/2}}
 \geq c\frac{n^{13/4}}{m^{5/2}}.
\end{equation}
Notice that a union bound over all $n!$ targets was not used.

\section{The complete nonlinear remainder on the binary family}
Write $R_\sigma-1=S_\sigma+U_\sigma$ and put
\[
 A=\frac{n^{13/4}}{m^{5/2}},\qquad t=m/n.
\]
We record explicitly the strengthened remainder estimate extracted from
the companion proof \emph{A refined uniform bound for random palindrome
blocks}.
For $p=4n$, $m\geq C_0 n^{13/10}$ with a sufficiently large absolute
$C_0$, and sufficiently large $n$,
\begin{equation}\label{eq:remainder}
 \|U_\sigma\|_p\leq C_1 A(A+t^{-1}).
\end{equation}
Here $C_1$ does not depend on $C_0,n,m$. To verify this, split the exact
expansion by its total selected segment length $s$ as in that proof.
For $s\leq n/2$, its colored coefficient estimate gives, for each
Hoeffding level $\ell$,
\[
 \|T_\ell^{\mathrm{short}}\|_p\leq a^\ell,
 \qquad a=2\sqrt{160}\,n^{11/4}\sqrt p/m^{5/2}=C_2 A.
\]
The levels $\ell\geq2$ contribute at most $a^2/(1-a)$.
In level one, remove segments of length three, which give exactly
$S_\sigma$ for $n\geq6$. The remaining size sum, with
$z=n^2/(2m^2)\leq1/2$, satisfies
\[
 \sum_{k\geq4}kz^k
 =\frac{z^4(4-3z)}{(1-z)^2}\leq10z^4
 \leq\frac{n^8}{m^8}.
\]
This gives the level-one remainder bound $a/t$.
For $s>n/2$, the same tilted generating-function estimate is
\[
 \|T^{\mathrm{long}}\|_p
 \leq\sqrt2\exp\left(-\frac n{12}\log t+\frac{4pn}{t^4}\right).
\]
Since $p=4n$ and $t\geq C_0 n^{3/10}$, for sufficiently large $n$ this is
at most $\sqrt2t^{-n/24}\leq A/t$.
Choosing $C_0$ so that $a\leq1/2$ proves \eqref{eq:remainder}.

Markov's inequality and $|\mathcal F|\leq2^n$ show that, except on an
event of probability at most $2^ne^{-4n}$,
\[
 \max_{\sigma\in\mathcal F}|U_\sigma|
 \leq eC_1 A(A+t^{-1}).
\]
Choose the absolute constant $C_0$ large enough that this is less than
half the lower bound in \eqref{eq:triplelower}; indeed
$A\leq C_0^{-5/2}$ and $t^{-1}\leq C_0^{-1}$.
Combining the two probability estimates proves Theorem~\ref{thm:lower}.

\begin{corollary}\label{cor:threshold}
There is an absolute $\varepsilon_0>0$ such that if
$m=m(n)=o(n^{13/10})$, then the probability that the random palindrome
construction is pointwise relatively $\varepsilon_0$-fair tends to zero.
\end{corollary}
\begin{proof}
If $m<n/(4e)$, then for sufficiently large $n$, $(2m)^n<n!$; some output
order has probability zero, so its relative error is one.
Otherwise $m\to\infty$. Choose
$k=\lfloor(m/(2C))^{10/13}\rfloor$, where $C$ is the threshold constant
in Theorem~\ref{thm:lower}. Then $k\to\infty$, $k\leq n$ eventually, and
$m\geq Ck^{13/10}$. Restrict the construction to any fixed $k$ labels.
The induced blocks are again independent uniform random palindromes,
and $k^{13/4}/m^{5/2}$ is bounded below by a positive absolute constant.
The theorem therefore gives a fixed positive lower bound for the maximal
relative error of this restriction, with probability tending to one.
Every relative error of a marginal order is an average of the full-order
relative errors. Hence the full construction has at least as large a
maximum relative error. Choose $\varepsilon_0$ below the obtained constant.
\end{proof}

\section{An exact covariance identity for the triple process}
For one block, zero-extend any finitely supported weights $(w_a)$ to all
integers. The centered valley variables have covariances
\[
 \operatorname{Cov}(\delta_a,\delta_{a+j})=
 \begin{cases}1/2,&j=0,\\-1/4,&|j|=1,\\1/20,&|j|=2,\\0,&|j|\geq3.\end{cases}
\]
Adjacent valleys cannot occur together. For distance two the joint
probability is
$\int_0^1(x-x^2/2)^2\,dx=2/15$, which proves the displayed values.
If $\Delta w_a=w_{a+1}-w_a$, direct expansion gives
\[
 \Var\Bigl(\sum_aw_a\delta_a\Bigr)
 =\frac1{10}\|w\|_2^2+\frac1{20}\|\Delta w\|_2^2
                  +\frac1{20}\|\Delta^2w\|_2^2.
\]
In particular, the cancellations between neighboring valleys do not remove
the variance of a slowly varying weighted sum.
For binomial weights of degree $d$, the exact integral identity
\[
 \int_0^1\sum_{a=0}^d\binom da^2p^{2a}(1-p)^{2d-2a}\,dp
 =\frac{4^d}{(2d+1)\binom{2d}d}
 \sim\frac{\sqrt\pi}{2\sqrt d}
\]
follows by the beta integral and the convolution identity
$\sum_a\binom{2a}a\binom{2d-2a}{d-a}=4^d$.

\paragraph{Dependency.}
The remainder estimate uses the complete coefficient-level inequalities,
not merely the final theorem, in
\texttt{refined\_random\_palindrome\_bound.tex}. That companion proof cites
the primary hypercontractivity source and was independently reviewed by
two other agents. No claim of a lower bound for general approximately fair
dice is intended here.
\end{document}
